\section{\texorpdfstring{Values of the $L$-Functions of $E_{W''}$}{Values of the L-Functions of E\_W''}}
\label{lv:section}

By \eqref{dh:YW}, $Y_{W''}(\sigma)$ is determined by $Y_E(\sigma)$ and by the
values at $\sigma$ of $576$ further $L$-functions: $448$ of degree $4$ over
$\Q$, the twists $L(s,\rho_o\otimes\chi_w)$ of the seven dihedral
representations, and $128$ of degree $8$, the functions
$L(s,\rho_a\otimes\rho_b\otimes\chi_w)$. They are needed at the $21$
abscissae of Table~\ref{ce:kernel-table}, all in $(1,101/100)$; the values
were computed at $36$ abscissae, from which the linear programme of
Section~\ref{ce:section} chose these $21$. Proposition~\ref{an:genus-value}
gives $Y_E(\sigma)$. This section encloses the $576$ values
(Proposition~\ref{lv:values}); Lemma~\ref{ce:values} uses them.

The $192$ twists of the forms $\psi_{10},\psi_{23},\psi_{24}$ were evaluated
with the certified approximate functional equation for the gamma factor
$\Gamma_\C(s)^2$ of \cite[Section~4.6, Lemma~4.4]{Naslund0418235}, with the
tail bounds of its Section~4.7 and the kernel enclosures of its
Section~4.9, which is not restated here. Its supplementary archive
implements it in the program \texttt{nonpositive-afe.py} (kernel type
\texttt{quartic}), which certifies enclosures of the kernels, of the
interpolation errors and of the truncation errors; the program
\texttt{dafe.py} applies it to these $192$ functions. The kernels of
Subsection~\ref{lv:afe-general} evaluate the same $192$ functions
independently (Subsection~\ref{lv:checks}, item~(3)), with radius about
$10^{-13}$; a reader may use those values instead, which only enlarges the
margin. The other $384$ functions were evaluated with the kernels of
Subsection~\ref{lv:afe-general}, which cover the four gamma factors of
Lemma~\ref{dh:gamma}.

\subsection{An approximate functional equation for products of gamma
functions}\label{lv:afe-general}
This subsection proves an approximate functional equation
(Lemma~\ref{lv:afe}) and the properties of its kernels
(Lemma~\ref{lv:kernel}) for the four gamma factors of Lemma~\ref{dh:gamma};
Subsection~\ref{lv:computation} uses them. In this section $d\in\{1,2\}$ is
a parameter of the gamma factor, not the degree $[F:\Q]$ of
Section~\ref{an:section}.

Let $d\in\{1,2\}$, let $m\ge1$ and $\mu_1,\dots,\mu_m\in\{0,\frac12\}$ with at
least one $\mu_j=0$, and let $C_G,\kappa_G>0$. The gamma factors of
Lemma~\ref{dh:gamma} have the form
$G(s)=C_G\kappa_G^{-s}\prod_{j=1}^m\Gamma(s/d+\mu_j)$:
\[
 \begin{aligned}
 \Gamma_\C(s)^4&=16(2\pi)^{-4s}\Gamma(s)^4,&
 \Gamma_\C(s)^2&=4(2\pi)^{-2s}\Gamma(s)^2,\\
 \Gamma_\R(s)^3\Gamma_\R(s+1)&=\pi^{-1/2}\pi^{-2s}\Gamma(\tfrac s2)^3\Gamma(\tfrac s2+\tfrac12),&
 \Gamma_\R(s)\Gamma_\R(s+1)^3&=\pi^{-3/2}\pi^{-2s}\Gamma(\tfrac s2)\Gamma(\tfrac s2+\tfrac12)^3,
 \end{aligned}
\]
with $(d,\kappa_G)=(1,(2\pi)^4)$, $(1,(2\pi)^2)$, $(2,\pi^2)$ and $(2,\pi^2)$.
Define the kernels
\begin{gather*}
 h(x)=\frac1{2\pi i}\int_{(\eta)}\prod_j\Gamma(u+\mu_j)\,x^{-u}\,du\qquad(x>0,\ \eta>0),\\
 \mathcal I_b(x)=\int_1^\infty r^{b-1}h(xr)\,dr\qquad(b\in\R).
\end{gather*}
Thus $h$ is the inverse Mellin transform of $\prod_j\Gamma(u+\mu_j)$, the
multiplicative convolution of the functions $x^{\mu_j}e^{-x}$; it is positive
and decays exponentially, so $\mathcal I_b(x)$ converges for every real $b$.

\begin{lemma}\label{lv:afe}
Let $L(s)=\sum_{n\ge1}a_nn^{-s}$ with real coefficients of polynomial growth,
let $Q>0$, and suppose that $\Lambda(s)=Q^{s/2}G(s)L(s)$ extends to an
entire function of finite order with $\Lambda(1-s)=\Lambda(s)$. Put
$t=\kappa_G/\sqrt Q$ and $x_n=(tn)^d$. Then for real $s>1$
\[
 L(s)=\frac{t^s}{\prod_j\Gamma(s/d+\mu_j)}\sum_{n\ge1}a_n
 \bigl(\mathcal I_{s/d}(x_n)+\mathcal I_{(1-s)/d}(x_n)\bigr),
\]
and the series converges absolutely.
\end{lemma}

\begin{proof}
For $y>0$ put $\Phi(y)=\sum_na_nh(x_ny)$. Since
$\int_0^\infty h(x)x^{u-1}dx=\prod_j\Gamma(u+\mu_j)$ for $\operatorname{Re}u>0$
and $Q^{s/2}\kappa_G^{-s}n^{-s}=x_n^{-s/d}$, termwise integration gives, for
$\operatorname{Re}(du)>1$,
\[
 \int_0^\infty\Phi(y)y^{u-1}dy=\prod_j\Gamma(u+\mu_j)\sum_na_nx_n^{-u}
 =\Lambda(du)/C_G.
\]
The right side is entire and invariant under $u\mapsto1/d-u$. As in the
proof of Lemma~\ref{an:balanced-afe}, $L$ has polynomial growth in vertical
strips by the Phragm\'en--Lindel\"of principle, so by Stirling's formula
$\Lambda(du)$ decays exponentially in $|\operatorname{Im}u|$ on vertical
strips. Mellin inversion on a line $\operatorname{Re}u=c>1/d$ and a shift to
$\operatorname{Re}u=1/d-c$, where no poles occur, give
$\Phi(y)=y^{-1/d}\Phi(1/y)$. Splitting the Mellin integral at $y=1$ and
substituting $y\mapsto1/y$ below $1$,
\begin{align*}
 \frac{\Lambda(s)}{C_G}&=\int_1^\infty\Phi(y)\bigl(y^{s/d-1}+y^{(1-s)/d-1}\bigr)dy\\
 &=\sum_na_n\bigl(\mathcal I_{s/d}(x_n)+\mathcal I_{(1-s)/d}(x_n)\bigr),
\end{align*}
by termwise integration, which the exponential decay of $h$ justifies.
Finally
\[
 L(s)=\frac{\Lambda(s)}{C_G\,t^{-s}\prod_j\Gamma(s/d+\mu_j)}.
\]
\end{proof}

\begin{lemma}\label{lv:kernel}
\begin{enumerate}
\item[(a)] $h$ and every $\mathcal I_b$ are completely monotone on
 $(0,\infty)$; in particular $|\mathcal I_b^{(k)}|$ is nonincreasing for every
 $k\ge0$.
\item[(b)] $x\mathcal I_b'(x)+b\mathcal I_b(x)=-h(x)$.
\item[(c)] $h(x)$, and $x^b\mathcal I_b(x)$ when no pole of
 $\prod_j\Gamma(b+\mu_j+u)$ lies at $u=0$, are given by convergent series of
 residues at the poles of the gamma factors, at $u=-n-\mu$ for $h$ and at
 $u=-n-\mu-b$ for $x^b\mathcal I_b$ ($n\ge0$, $\mu\in\{\mu_j\}$), and
 $x^b\mathcal I_b(x)$ has the additional term $\prod_j\Gamma(b+\mu_j)$ from
 the pole of $1/u$. When $b<0$ the pole with $n=0$, $\mu=0$ lies at $u=-b>0$,
 and the line of integration is taken to its right.
\item[(d)] For $\eta>0$ with $b+\mu_j+\eta>0$ for all $j$,
 $|\mathcal I_b(x)|\le C_{b,\eta}\,x^{-b-\eta}$, where
 $C_{b,\eta}=\prod_j\Gamma(b+\mu_j+\eta)\cdot\frac1{2\pi}\int_\R
 \prod_j\bigl(1+\tfrac{y^2}{(b+\mu_j+\eta)^2}\bigr)^{-1/2}\frac{dy}{|\eta+iy|}$.
\end{enumerate}
\end{lemma}

\begin{proof}
(a) If $\mu_{j_0}=0$, then $h(x)=\int_0^\infty e^{-x/v}\varrho(v)\,dv/v$, where
$\varrho\ge0$ is the convolution of the other factors $x^{\mu_j}e^{-x}$; this
is a positive combination of the completely monotone functions $e^{-x/v}$.
Then $\mathcal I_b(x)=\int_1^\infty r^{b-1}h(xr)\,dr$ is also one, and the
derivatives of a completely monotone function alternate in sign and have
nonincreasing absolute values.
(b) Substituting $v=xr$, $\mathcal I_b(x)=x^{-b}\int_x^\infty v^{b-1}h(v)\,dv$;
now differentiate.
(c) The Mellin transform of $x^b\mathcal I_b(x)=\int_x^\infty v^{b-1}h(v)dv$ is
$\prod_j\Gamma(b+\mu_j+u)/u$; shift the inversion integral to the left. The
residues decay faster than any geometric sequence.
(d) Write $x^b\mathcal I_b(x)=\frac1{2\pi i}\int_{(\eta)}\prod_j
\Gamma(b+\mu_j+u)x^{-u}du/u$ and use
$|\Gamma(a+iy)|\le\Gamma(a)(1+y^2/a^2)^{-1/2}$ for $a>0$, which follows from
the product formula for $|\Gamma(a+iy)|^2$.
\end{proof}

\subsection{The computation}\label{lv:computation}
This subsection describes how the kernels and the coefficients are computed
and how their errors are bounded (Lemmas~\ref{lv:remainders}
and~\ref{lv:coefficients}); the resulting enclosures,
Proposition~\ref{lv:values}, are the input of Lemma~\ref{ce:values}. For each
$L$-function let $a_n$ be its coefficients and $N$ the truncation point of
the series of Lemma~\ref{lv:afe} (the column ``terms'' of
Table~\ref{lv:results}).

For each gamma factor, the supplementary program \texttt{deg8/gkernel.py}
encloses $\mathcal I_{\sigma/d}$ and $\mathcal I_{(1-\sigma)/d}$ in Arb on a
grid of cells: the points $\xi_0>\xi_1>\dots$ are $\xi_0=6144$ for
$\Gamma_\C(s)^4$, $512$ for $\Gamma_\C(s)^2$ and $2^{14}$ for the mixed
factors, and $\xi_{i+1}=31\xi_i/32$ rounded down to $53$ significant bits, down to
$2^{-40}$. At each $\xi_i$ the Taylor coefficients of order $0$ to $20$ of
$\mathcal I_b$ are computed from the residue series of
Lemma~\ref{lv:kernel}(c), with the working precision raised to absorb the
cancellation in these alternating series, whose terms reach about
$\exp(m x^{1/m})$, and with the truncated tail bounded by Cauchy's estimate on
circles of radius $1/2$ or $1/4$ around the poles; the recursion of
Lemma~\ref{lv:kernel}(b) gives the higher coefficients from those of $h$. The
next lemma bounds the error on each cell and beyond the last cell used.

\begin{lemma}\label{lv:remainders}
Let $b\in\R$.
\begin{enumerate}
\item[(a)] For $i\ge0$ and $x\in[\xi_{i+1},\xi_i]$,
 \[
  \Bigl|\mathcal I_b(x)-\sum_{k=0}^{20}\frac{\mathcal I_b^{(k)}(\xi_i)}{k!}
  (x-\xi_i)^k\Bigr|\le\frac{|\mathcal I_b^{(21)}(\xi_{i+1})|}{21!}
  (\xi_i-\xi_{i+1})^{21}.
 \]
\item[(b)] Let $\eta$ and $C_{b,\eta}$ be as in Lemma~\ref{lv:kernel}(d), let
 $0<s_0<d(b+\eta)$, let $t>0$ and $N\ge1$, and let $x_n=(tn)^d$. Then for
 real numbers $a_n$
 \[
  \sum_{n>N}|a_n|\,|\mathcal I_b(x_n)|\le C_{b,\eta}\,t^{-d(b+\eta)}
  N^{s_0-d(b+\eta)}\sum_{n\ge1}|a_n|n^{-s_0}.
 \]
\end{enumerate}
\end{lemma}

\begin{proof}
(a) By Taylor's theorem the difference is
$\mathcal I_b^{(21)}(\vartheta)(x-\xi_i)^{21}/21!$ for some $\vartheta\in[x,\xi_i]$. By
Lemma~\ref{lv:kernel}(a), $|\mathcal I_b^{(21)}(\vartheta)|\le|\mathcal
I_b^{(21)}(\xi_{i+1})|$, since $\vartheta\ge \xi_{i+1}$; and $|x-\xi_i|\le \xi_i-\xi_{i+1}$.
(b) This is Rankin's trick. By Lemma~\ref{lv:kernel}(d),
$|\mathcal I_b(x_n)|\le C_{b,\eta}(tn)^{-d(b+\eta)}$, and for $n>N$,
$n^{-d(b+\eta)}=n^{-s_0}n^{s_0-d(b+\eta)}\le n^{-s_0}N^{s_0-d(b+\eta)}$.
\end{proof}

The coefficient programs \texttt{deg8/dcoef4.c} (degree $4$) and
\texttt{deg8/octcoef.c} (degree $8$) compute, for each $L$-function and each
cell $i$ (the $n$ with $x_n\in(\xi_{i+1},\xi_i]$), the moments
$S_{ik}=\sum_na_nu_n^k$, $u_n=x_n/\xi_i-1$, for $k\le20$, and
$A_i=\sum_n|a_n|$, in double precision with compensated (Kahan) summation.
These moments differ from those of Subsection~\ref{an:kernel-enclosures}:
the variable is $u_n$, and there is no weight $n^\nu$. The program
\texttt{deg8/leval.py} turns kernels and moments into enclosures
\[
 L(\sigma)\in\frac{t^\sigma}{\prod_j\Gamma(\sigma/d+\mu_j)}\Bigl(\sum_i\sum_kp_{ik}S_{ik}
 \pm\bigl(\textstyle\sum_iA_iR_i+\text{tail}\bigr)\Bigr),
\]
where $p_{ik}$ are the Taylor coefficients of the kernels at $\xi_i$ in the
variable $u$, $R_i$ the cell remainders of Lemma~\ref{lv:remainders}(a) and
``tail'' the bound of Lemma~\ref{lv:remainders}(b), for the two kernels
$\mathcal I_{\sigma/d}$ and $\mathcal I_{(1-\sigma)/d}$. It treats each moment
as a ball whose radius bounds the floating-point errors: the error of $t$ in
double precision, bounded in Arb; the errors of $x_n$, $u_n$ and the powers,
a few units in the last place each ($|u_n|\le1-\xi_{i+1}/\xi_i$ exactly, since the
cells are cut at exact thresholds); and the error of Kahan summation, at most
$(2\cdot2^{-53}+O(n2^{-106}))$ times the sum of the absolute values of the
terms \cite[Section~4.3]{Higham2002}, plus the merging of the thread
accumulators. For the degree-$4$ functions the Rankin tail uses
$|a_n|\le\tau_4(n)$, where $\tau_4(n)$ is the number of ordered
factorizations of $n$ into four positive integers, and
$\sum\tau_4(n)n^{-s_0}=\zeta(s_0)^4$. For degree $8$ it uses the Euler
product of $\sum|a_n|n^{-s_0}$: PARI's Euler factors for $p\le p_0$, where
$p_0=2^{18}$ for the pairs $(19,23)$, $(19,24)$ and $2^{19}$ for the other
six; the sums computed by \texttt{octcoef.c} for the primes in $(p_0,N]$; the
bound $\binom{k+7}7$ for the coefficients of prime powers beyond; and the
bound $\pi(x)\le1.25506\,x/\log x$ of Rosser and Schoenfeld
\cite{RosserSchoenfeld1962} for the primes beyond $N$.

The coefficients at large primes follow from the next lemma. Recall from
Lemma~\ref{dh:Lfunctions} the elements $\gamma_o\in B_{c_o}$ and the fields
$M_{o,w}$, $B_{ab}$ and $B_{ab,w}$.

\begin{lemma}\label{lv:coefficients}
\begin{enumerate}
\item[(a)] Let $\psi_o$ be a form of Table~\ref{dh:W-table}, let
 $w\in\Ftwo^8$, and let $p\notin\{2,3,5,241\}$ be prime. The coefficient
 $a_p$ of $L(s,\rho_o\otimes\chi_w)$ is $0$ if $p$ is inert in $B$. If $p$
 splits in $B$, then
 \[
  a_p=\sum_{\mathfrak P}\Bigl(\frac{\gamma_o\alpha_w}{\mathfrak P}\Bigr),
 \]
 the sum of quadratic residue symbols over the primes $\mathfrak P$ of
 $B_{c_o}$ of norm $p$, provided $\gamma_o$ is a unit at these primes.
\item[(b)] Let $(a,b)$ be a pair of Lemma~\ref{dh:W-structure}, let
 $w\in\Ftwo^8$, and let $\mathfrak q\notin S$ be a prime of $B$ of degree
 one and norm $q\ne241$, with Frobenius $g\in G_{W''}$. For $o\in\{a,b\}$: if
 the image of $g$ in $\Gal(M_o/B)$ is central, then $\mathfrak q$ splits completely in
 $F_{0,o}$; let $\epsilon_o(\mathfrak q)\in\{\pm1\}$ be the quadratic residue
 symbol of $\gamma_o$ modulo a prime of $B_{c_o}$ above $\mathfrak q$,
 provided $\gamma_o$ is a unit there. Then
 $\tr(\rho_a\otimes\rho_b\otimes\chi_w)(g)=0$ unless the image of $g$ is
 central in both $\Gal(M_a/B)$ and $\Gal(M_b/B)$, and then it equals
 $4\epsilon_a(\mathfrak q)\epsilon_b(\mathfrak q)\chi_w(\mathfrak q)$. The
 coefficient $a_q$ of $L(s,\rho_a\otimes\rho_b\otimes\chi_w)$ is the sum of
 these traces over the two primes of $B$ above $q$.
\end{enumerate}
\end{lemma}

\begin{proof}
The fields $M_{o,w}$ and $B_{ab,w}$ are subfields of $\EW$ (proof of
Lemma~\ref{dh:Lfunctions}), which is unramified over $B$ outside $S$ since
$\EW\subseteq K$; so $p$, $q$ and $\mathfrak q$ are unramified in them.

(a) By Lemma~\ref{dh:Lfunctions}(a), $L(s,\rho_o\otimes\chi_w)$ is the
$L$-function of the quadratic character $\chi$ of $B_{c_o}$ that cuts out
$M_{o,w}=B_{c_o}(\sqrt{\gamma_o\alpha_w})$, so $a_p$ is the sum of
$\chi(\mathfrak P)$ over the primes $\mathfrak P$ of $B_{c_o}$ of norm $p$.
If $p$ is inert in $B$, there are none. Otherwise $\mathfrak P$ is odd and
unramified in $M_{o,w}$, and $\gamma_o\alpha_w$ is a unit at $\mathfrak P$,
since $\alpha_w$ is a unit outside $S$; so $\chi(\mathfrak P)=1$ exactly when
$\gamma_o\alpha_w$ is a square modulo $\mathfrak P$.

(b) The trace of a tensor product is the product of the traces, and
$\chi_w(g)=\chi_w(\mathfrak q)$. The character of the two-dimensional
irreducible representation of $\mathrm{Dih}_4$ is $2$ at $1$, $-2$ at the
central element $z$, and $0$ elsewhere. The image of $g$ in $\Gal(M_o/B)$ is
central exactly when it lies in $\langle z\rangle=\Gal(M_o/F_{0,o})$, that
is, when $\mathfrak q$ splits completely in $F_{0,o}$; then the primes of
$B_{c_o}$ above $\mathfrak q$ have degree one, and the image is $1$ exactly
when it fixes $\sqrt{\gamma_o}=(1+u_o)\sqrt{\beta_o}$, since $z$ negates
it, that is, when $\gamma_o$ is a square modulo a prime of $B_{c_o}$ above
$\mathfrak q$. So $\tr\rho_o(g)=2\epsilon_o(\mathfrak q)$ in that case, and
$0$ otherwise. Finally, the coefficient at $q$ of an Artin $L$-function over
$B$ is the sum of the traces of the Frobenius elements of the primes of $B$
of norm $q$.
\end{proof}

\emph{Coefficients of degree $4$.} For $p\le3000$ the Euler factors are
PARI's ($\zeta_{M_{o,w},p}/\zeta_{B_{c_o},p}$, program \texttt{deg8/exportg.gp}).
For $p>3000$ only $a_p$ enters, since the truncation point $N$ is at most
$6.7\cdot10^6<3000^2$, and only primes $p$ that split in $B$; then $a_p$ is
given by Lemma~\ref{lv:coefficients}(a). The program \texttt{dcoef4.c} stops
if a residue symbol vanishes; this does not occur.

\emph{Coefficients of degree $8$.} For $p\le p_0$ the Euler factors are PARI's
($\zeta_{B_{ab,w},p}/\zeta_{B_{ab},p}$, program \texttt{deg8/export8.gp}).
The truncation point is $N=5.28\cdot10^{10}$ for the first two pairs and
$1.055\cdot10^{11}$ for the others, which is less than $p_0^2$. So every
$n\le N$ is either $p_0$-smooth or $n_1q$ with a prime $q>p_0$ and $n_1<p_0$,
and for $q>p_0$ only $a_q$ enters, and only primes $q$ that split in $B$.
The program \texttt{octcoef.c} sieves the primes $q\in(p_0,N]$
($2.23\cdot10^9$ of them for the first two pairs, $4.34\cdot10^9$ for the
others), evaluates the rule of Lemma~\ref{lv:coefficients}(b), and enumerates
the $n=n_1q$ and the $p_0$-smooth $n$ with $a_n\ne0$. A prime $\mathfrak q$
at which the images of the Frobenius are central but $\gamma_o$ is not a unit
at a prime of $B_{c_o}$ above $\mathfrak q$ would be listed as exceptional;
none occurs. Before each run the program compares the rule with PARI's Euler
factors at every unramified prime in $(1000,p_0]$, for all sixteen twists of
the pair: these are $22832$ primes for the first two pairs and $43222$ for the
others, with between $1395$ and $2722$ primes of $B$ per pair at which both
quotients are central.

\subsection{Checks}\label{lv:checks}
The following checks are independent of the coefficient rules and of the
kernels, or of part of them; none is used in the proof.
\begin{enumerate}
\item[(1)] \emph{Euler products.} At $s=19/10$ the Euler products of the
 degree-$8$ functions converge absolutely. For every pair, the program
 \texttt{deg8/octtest.py} compares the value given by the approximate
 functional equation (radius about $10^{-14}$) with the Euler product
 (radius about $3\cdot10^{-10}$): all $128$ functions agree, with midpoints
 within $8\cdot10^{-16}$. A wrong conductor, gamma factor, root number or
 coefficient rule would show here.
\item[(2)] \emph{Kernels.} For each of the four gamma factors, \texttt{gkernel.py}
 reproduces PARI's \texttt{lfun} values of products of quadratic Dirichlet
 $L$-functions with that gamma factor, within the certified radii.
\item[(3)] \emph{Two implementations.} \texttt{dcoef4.c} with
 \texttt{leval.py} reproduces the values computed with the archive kernels of
 \cite{Naslund0418235} by \texttt{dafe.py}: for the twists of
 $\psi_{19}$ at $801/800$, and for all $6912$ values of the $192$ twists of
 $\psi_{10},\psi_{23},\psi_{24}$ at the $36$ abscissae, where all balls overlap
 and the midpoints agree to $3.4\cdot10^{-16}$.
\item[(4)] \emph{Coefficients.} The moments of \texttt{dcoef4.c} and
 \texttt{octcoef.c} agree with independent Python implementations up to
 $10^6$ and $2\cdot10^6$; the coefficients of the $192$ twists of
 $\psi_{10},\psi_{23},\psi_{24}$ agree with PARI's Dirichlet series
 $\zeta_{M_{o,w}}/\zeta_{B_{c_o}}$ up to $20000$, and for sample twists PARI's
 \texttt{lfun} with the same data satisfies the functional equation to
 $2^{-41}$ or better.
\item[(5)] \emph{Reviews.} Independent reviews reproduced the Euler-factor
 identity of Lemma~\ref{dh:EW}(e) at all primes up to $1000$ and at $2,3,5$;
 the degree-$8$ Euler factors at $57$ primes of vector $0$; all degree-$8$
 conductors by a different construction; and PARI's \texttt{lfun} values of
 the twists of $\psi_{20}$, $\psi_{17}$ and $\psi_7$ inside the stored
 enclosures. No reviewer recomputed the degree-$8$ values, whose conductors
 are beyond the reach of general-purpose software.
\end{enumerate}

\subsection{Results}\label{lv:results-section}
This subsection states the enclosures; with \eqref{dh:YW} they are used in
Lemma~\ref{ce:values}.

\begin{proposition}\label{lv:values}
For each abscissa $\sigma_j$ of Table~\ref{ce:kernel-table} and each of the
$576$ $L$-functions $L$ in \eqref{dh:YW}, the directories \texttt{hp} and
\texttt{hp3} contain an interval that contains $L(\sigma_j)$, of radius at
most the entry of Table~\ref{lv:results} (relative to the value, or absolute
for the archive evaluation). With Proposition~\ref{an:genus-value} and
\eqref{dh:YW} they give an enclosure of $Y_{W''}(\sigma_j)$; for instance
$Y_{W''}(1001/1000)=0.050492599$ with radius $3.4\cdot10^{-10}$.
\end{proposition}

\begin{proof}[Proof (computer-assisted)]
By Lemmas~\ref{dh:Lfunctions}(c) and~\ref{dh:gamma}, each of the $576$
functions satisfies the hypotheses of Lemma~\ref{lv:afe}, with $Q$ its
conductor and $G$ its gamma factor. For the $384$ functions of the forms
$\psi_{19},\psi_{20},\psi_{17},\psi_7$ and of the eight pairs, the program
\texttt{deg8/leval.py} evaluates the series of Lemma~\ref{lv:afe} as
described in Subsection~\ref{lv:computation}: with the kernels of
\texttt{deg8/gkernel.py} (Lemma~\ref{lv:kernel}), the coefficients from
PARI's Euler factors and Lemma~\ref{lv:coefficients}, the cell remainders and
tails of Lemma~\ref{lv:remainders}, and the floating-point errors bounded as
described there. The $192$ twists of $\psi_{10},\psi_{23},\psi_{24}$ are
evaluated by \texttt{dafe.py} with the approximate functional equation of
\cite[Lemma~4.4]{Naslund0418235}, or alternatively, with radius about
$10^{-13}$, by the same pipeline as the other functions
(Subsection~\ref{lv:checks}, item~(3)).
\end{proof}

\begin{table}[htbp]
\centering\small
\caption{The $L$-functions in \eqref{dh:YW}: number, gamma factor, terms of
the approximate functional equation, and the largest radius of the
enclosures of their values at the abscissae (relative, except for $Y_E$ and
for the archive evaluation, where it is absolute). All data are in the
directories \texttt{hp} and \texttt{hp3}.}
\label{lv:results}
\begin{tabular}{lrlll}
\toprule
functions & number & gamma factor & terms & radius\\
\midrule
quadratic, $E$ & $255$ & $\Gamma_\R(s+\nu)^2$ or $\Gamma_\C(s)$ & $8\lfloor\sqrt{Q_e}\rfloor+1$ & $1.3\cdot10^{-19}$ (in $Y_E$)\\
$\psi_{10},\psi_{23},\psi_{24}$ & $192$ & $\Gamma_\C(s)^2$ & $999999$ & $6.6\cdot10^{-9}$\\
$\psi_{19},\psi_{20}$ & $128$ & $\Gamma_\C(s)^2$, mixed & up to $6.7\cdot10^6$ & $6.4\cdot10^{-13}$\\
$\psi_{17},\psi_7$ & $128$ & $\Gamma_\C(s)^2$ & up to $6.7\cdot10^6$ & $2.2\cdot10^{-13}$\\
pairs $(19,23),(19,24)$ & $32$ & $\Gamma_\C(s)^4$ & up to $5.28\cdot10^{10}$ & $5.1\cdot10^{-10}$\\
other six pairs & $96$ & $\Gamma_\C(s)^4$ & up to $1.055\cdot10^{11}$ & $5.1\cdot10^{-10}$\\
\bottomrule
\end{tabular}
\end{table}
